\documentclass[11pt,a4paper]{article}

% ============================================================
% Maliyadeva Simulation and Modelling Society
% Constitution
% ============================================================

% ------------------------------------------------------------
% Encoding and Typography
% ------------------------------------------------------------
\usepackage[T1]{fontenc}
\usepackage[utf8]{inputenc}
\usepackage{lmodern}
\usepackage{microtype}

% ------------------------------------------------------------
% Page Layout
% ------------------------------------------------------------
\usepackage[
    a4paper,
    top=1in,
    bottom=1in,
    left=1.05in,
    right=1.05in,
    headheight=14pt
]{geometry}

% ------------------------------------------------------------
% Language and Text
% ------------------------------------------------------------
\usepackage[english]{babel}
\usepackage{csquotes}

% ------------------------------------------------------------
% Lists
% ------------------------------------------------------------
\usepackage{enumitem}

\setlist[enumerate]{
    leftmargin=2em,
    itemsep=0.25em,
    topsep=0.5em,
    parsep=0pt
}

% ------------------------------------------------------------
% Section Formatting
% ------------------------------------------------------------
\usepackage{titlesec}

% Articles
\titleformat{\section}
    {\normalfont\bfseries\centering}
    {}
    {0pt}
    {\MakeUppercase}

\titlespacing*{\section}
    {0pt}
    {2.5em}
    {1.25em}

% Sections within Articles
\titleformat{\subsection}
    {\normalfont\bfseries}
    {}
    {0pt}
    {}

\titlespacing*{\subsection}
    {0pt}
    {1.5em}
    {0.5em}

% ------------------------------------------------------------
% Paragraph Formatting
% ------------------------------------------------------------
\setlength{\parindent}{0pt}
\setlength{\parskip}{0.65em}

% Prevent awkward paragraph breaks where possible
\widowpenalty=10000
\clubpenalty=10000

% ------------------------------------------------------------
% Headers and Footers
% ------------------------------------------------------------
\usepackage{fancyhdr}

\pagestyle{fancy}

\fancyhf{}

\fancyhead[L]{\small\textit{Maliyadeva Simulation and Modelling Society}}
\fancyhead[R]{\small\textit{Constitution}}

\fancyfoot[C]{\small\thepage}

\renewcommand{\headrulewidth}{0.4pt}
\renewcommand{\footrulewidth}{0pt}

% ------------------------------------------------------------
% Hyperlinks and PDF Metadata
% ------------------------------------------------------------
\usepackage[
    hidelinks,
    pdftitle={Constitution of the Maliyadeva Simulation and Modelling Society},
    pdfauthor={Maliyadeva College, Kurunegala},
    pdfsubject={Constitution of the Maliyadeva Simulation and Modelling Society},
    pdfkeywords={
        Maliyadeva College,
        Simulation,
        Modelling,
        Society,
        Constitution,
        MSMS
    }
]{hyperref}

% ------------------------------------------------------------
% Better Cross-References
% ------------------------------------------------------------
\usepackage[nameinlink,noabbrev]{cleveref}

% ------------------------------------------------------------
% Document Metadata
% ------------------------------------------------------------
\title{
    \vspace{-1.5cm}
    {\Large\bfseries CONSTITUTION}\\[0.5em]
    {\large\bfseries OF THE}\\[0.5em]
    {\LARGE\bfseries MALIYADEVA SIMULATION AND MODELLING SOCIETY}\\[0.75em]
    {\normalsize Maliyadeva College, Kurunegala}
}

\author{}
\date{}

% ------------------------------------------------------------
% Custom Commands
% ------------------------------------------------------------

% Article reference
\newcommand{\article}[1]{\textbf{Article #1}}

% Section reference
\newcommand{\sectionref}[1]{\textbf{Section #1}}

% Formal quoted terms
\newcommand{\definedterm}[1]{\enquote{#1}}

% ------------------------------------------------------------
% Document
% ------------------------------------------------------------

\begin{document}

\maketitle
\thispagestyle{empty}

\clearpage


\section*{Article I --- General Provisions}

\subsection*{Section 1. Name}
The official name of this organisation shall be the Maliyadeva Simulation and Modelling Society, referred to in this Constitution as ``the Society'' or by the acronym MSMS.

\subsection*{Section 2. Status and Affiliation}
The Society shall operate as an official student organisation affiliated with and governed under the administrative oversight of Maliyadeva College.

\subsection*{Section 3. Purpose and Objectives}
The purpose of the Society is to advance understanding and practice of computer simulation, computational modelling, and applied systems thinking among the students of Maliyadeva College.

The Society shall pursue this purpose through:
\begin{enumerate}
    \item Learning --- promoting interest, technical skill, and academic understanding in simulation, modelling, and applied systems thinking;
    \item Application --- enabling students to apply mathematical and computational techniques to practical problems;
    \item Research and projects --- conducting collaborative, research-inspired projects and maintaining reproducible technical standards;
    \item Mentoring --- developing the skills of members and facilitating the transfer of knowledge between experienced and newer members;
    \item Competition and external participation --- representing the Society in approved competitions and programmes;
    \item Outreach --- promoting simulation and modelling within the College and the wider community; and
    \item Responsible practice --- maintaining academic and technical integrity in all Society work.
\end{enumerate}

\section*{Article II --- Governing Authority and Hierarchy}

\subsection*{Section 1. Supremacy of Governing Instruments}
The governing instruments of the Society shall have the following order of precedence:
\begin{enumerate}
    \item the policies, rules, and requirements of Maliyadeva College;
    \item this Constitution;
    \item the Bylaws;
    \item Regulations adopted under the Bylaws, including the Membership Regulations, Financial Regulations, the Code of Conduct, and the Transitional Regulations;
    \item resolutions of the General Assembly;
    \item resolutions and decisions of the Executive Committee; and
    \item decisions of individual officers, committees, project teams, and other subordinate bodies.
\end{enumerate}

Where a lower-ranking instrument conflicts with a higher-ranking instrument, the higher-ranking instrument shall prevail, and the conflicting provision shall be void to the extent of the conflict.

\subsection*{Section 2. College Authority}
The Society operates under the ultimate authority of Maliyadeva College. The College retains authority over the existence, activities, and compliance of the Society, and may require referral of any matter to the appropriate College authority, direct that an activity be suspended or modified, or impose any other requirement permitted by College policy or applicable law.

The autonomy of the Society under this Constitution and the Bylaws shall be exercised subject to the College's authority.

\subsection*{Section 3. Severability}
If any provision of this Constitution is found to be invalid or unenforceable, the remaining provisions shall remain in effect to the extent permitted by applicable requirements.

\section*{Article III --- Membership}

\subsection*{Section 1. Eligibility}
All students currently enrolled at Maliyadeva College shall be eligible to apply for membership in the Society, subject to the procedures prescribed in the Bylaws and the Membership Regulations.

\subsection*{Section 2. Categories}
The Society shall consist of three membership categories:
\begin{enumerate}
    \item Associate Membership --- open to all students who wish to participate in the activities, workshops, and projects of the Society. Associate Members shall not possess voting rights and shall not be eligible for election or nomination to the Executive Committee.
    \item Regular Membership --- granted to an Associate Member who has demonstrated active and meaningful participation in the Society and who satisfies the progression requirements prescribed in the Bylaws and the Membership Regulations. Regular Members possess full membership and voting rights and shall be eligible for nomination and election to the Executive Committee, subject to the Bylaws.
    \item Honorary Membership --- conferred upon an individual who has made a distinguished contribution to the Society, modelling and simulation, education, research, or a related field, or who has provided exceptional service or support to the Society. Honorary Membership shall be conferred by resolution of the Executive Committee. Honorary Members may participate in activities and may be invited to provide advice, mentorship, lectures, or other support, but shall not possess voting rights or be eligible for election to the Executive Committee solely by virtue of their honorary membership.
\end{enumerate}

\subsection*{Section 3. Fundamental Rights}
All members shall have the rights:
\begin{enumerate}[leftmargin=*]
    \item Participation --- to participate in Society activities, workshops, projects, competitions, and other opportunities, subject to applicable requirements.
    \item Access to information --- to receive relevant information concerning Society activities and decisions affecting them.
    \item Expression --- to express opinions, propose ideas, and contribute to discussions in accordance with the Society's procedures.
    \item Fair treatment --- to be treated fairly and without unreasonable discrimination in matters concerning membership.
    \item Use of resources --- to access Society resources made available to members, subject to applicable rules.
    \item Due process --- where a member faces suspension, termination, or disciplinary action, to know the grounds and to respond in accordance with the applicable procedure.
\end{enumerate}
Associate Members shall have all rights applicable to members except voting rights and eligibility for election or nomination to the Executive Committee.

\subsection*{Section 4. Termination --- Principles}
Membership shall cease upon:
\begin{enumerate}
    \item Voluntary resignation --- a member may resign by submitting written notice to the Secretary. Resignation shall take effect upon receipt of the notice, unless a later effective date is specified.
    \item Loss of eligibility --- where a member no longer satisfies the eligibility requirements applicable to their membership category.
    \item Disciplinary removal --- where membership is terminated as a disciplinary sanction in accordance with the disciplinary procedures prescribed in the Bylaws and the Code of Conduct.
\end{enumerate}
The detailed procedures for the cessation of membership shall be prescribed in the Bylaws and the Membership Regulations.

\section*{Article IV --- Governing Bodies}

\subsection*{Section 1. General Assembly}
The General Assembly shall consist of all Regular Members in good standing and shall
serve as the principal deliberative and membership-governance body of
the Society.

The General Assembly shall:
\begin{enumerate}[leftmargin=*]
    \item elect the officers designated as elective offices under the
    Bylaws;
    \item approve amendments to this Constitution where required;
    \item hold the Executive Committee accountable for the affairs of
    the Society;
    \item remove elected officers in accordance with the Bylaws; and
    \item exercise such other powers as are assigned to it by this
    Constitution and the Bylaws.
\end{enumerate}

\subsection*{Section 2. Executive Committee}
The Executive Committee shall be the principal executive body of the
Society and shall manage its affairs in accordance with this
Constitution, the Bylaws, decisions of the General Assembly, and
applicable College requirements.

The offices, powers, duties, meetings, quorum, and procedures of the Executive
Committee shall be prescribed in the Bylaws.

\subsection*{Section 3. Relationship Between Governing Bodies}
The General Assembly shall hold the Executive Committee accountable for
the administration of the Society, and the Executive Committee shall
remain subject to this Constitution, the Bylaws, and decisions of the
General Assembly made within its authority.

Committees, project teams, and other subordinate bodies established
under the Bylaws shall exercise only the authority delegated to them and
shall remain accountable to the Executive Committee.

\subsection*{Section 4. Teacher-in-Charge}

The Society shall have a Teacher-in-Charge or Teachers appointed or
recognised by Maliyadeva College in accordance with applicable College
procedures.

The Teacher-in-Charge shall serve as the principal liaison between the
Society and the College and shall advise the Society on matters
concerning College requirements, student welfare, safety, finances, and
the proper conduct of Society activities.

The Teacher-in-Charge may require a matter to be referred to the
appropriate College authority, or require an activity to be suspended
or modified, where this is necessary to comply with College policy,
safety requirements, applicable law, or other requirements of the
College.

The Teacher-in-Charge may exercise such supervisory, approval,
countersignature, and adjudicative functions as are expressly assigned
to the office by this Constitution, the Bylaws, or applicable College
requirements, including functions relating to Society finances,
disciplinary matters, and the compliance of Society activities with
College requirements.

The Teacher-in-Charge shall not be an elected officer or member of the
Executive Committee solely by virtue of holding the position. Subject
to the authority of the College, the day-to-day governance of the
Society shall remain with the General Assembly and Executive Committee
in accordance with this Constitution and the Bylaws, except for the
specific functions expressly assigned to the Teacher-in-Charge under
this Constitution, the Bylaws, or applicable College requirements.

\section*{Article V --- Amendment of the Constitution}

\subsection*{Section 1. Proposal}
An amendment to this Constitution may be proposed:
\begin{enumerate}[leftmargin=*]
    \item by the Executive Committee acting by resolution; or
    \item by any Regular Member, provided that the proposed amendment is submitted in writing and supported by the signatures of at least ten Regular Members.
\end{enumerate}

\subsection*{Section 2. Approval}
An amendment to this Constitution shall require:
\begin{enumerate}[leftmargin=*]
    \item a quorum of at least two-thirds of all Regular Members; and
    \item the affirmative vote of at least two-thirds of the Regular Members present and voting.
\end{enumerate}
Where College approval is required for the amendment, that approval shall be obtained in addition to the approval of the General Assembly.

\subsection*{Section 3. Effective Date}
An amendment to this Constitution shall take effect upon approval by the General Assembly and upon any approval required by Maliyadeva College.

Unless expressly provided otherwise and permitted by applicable requirements, an amendment shall not operate retroactively or affect any act, decision, appointment, transaction, or other matter validly completed before its effective date.

\section*{Article VI --- Dissolution}

\subsection*{Section 1. Dissolution}
Subject to the authority of Maliyadeva College under Article II,
Section 2, the Society may be dissolved upon:
\begin{enumerate}[leftmargin=*]
    \item a resolution of the Executive Committee;
    \item a written petition supported by at least one-third of all
    Regular Members; or
    \item such other direction or requirement of Maliyadeva College as
    may apply.
\end{enumerate}

Where dissolution is initiated by the Society, it shall require the
affirmative vote of at least two-thirds of the Regular Members present
and voting at a General Assembly meeting at which at least two-thirds of
all Regular Members are present.

Dissolution shall take effect upon any approval or recognition required
by Maliyadeva College. Upon dissolution, the Society shall cease to
operate except as necessary to wind up its affairs.

\subsection*{Section 2. Assets, Records, and Obligations}
Upon dissolution, the Society's outstanding obligations shall be
settled in accordance with the Bylaws, the Financial Regulations, and
applicable College requirements.

Any remaining funds, property, records, intellectual property, and
digital or other assets of the Society shall be transferred, retained,
or otherwise dealt with as directed or approved by Maliyadeva College
and in accordance with applicable obligations.

No Society asset shall be distributed for the personal benefit of a
member or officer solely by reason of their membership or office.

\section*{Article VII --- Final Provisions}

\subsection*{Section 1. Effective Date}
This Constitution shall become effective upon its approval by the General Assembly in accordance with the applicable provisions of this Constitution and the Bylaws, and upon completion of any approval, recognition, or confirmation required by Maliyadeva College, unless a later effective date is expressly specified in the approval or by the College.

\subsection*{Section 2. Supersession and Repeal}
Upon the effective date of this Constitution, the Charter of the Simulation \& Modelling Society dated May 27, 2026, and all previous constitutions, charters, bylaws, regulations, and other internal governing documents of the Society are hereby repealed in their entirety.

\subsection*{Section 3. Transitional Provisions}
The Transitional Regulations of the Society shall apply from the effective date of this Constitution and shall remain in force only for so long as necessary to complete the transition to the governance framework established by this Constitution and the Bylaws.

\end{document}
